\section{第\ref{sec:dofa}章　\nt{DOPA}による学習}

シナプスの機構（パケット、同時検出器、カルシウム、可塑性の窓）と、予測誤差としてのドーパミンのふるまいは直接測定されている。規則が勾配に導くこと、ブロードキャストの代償は定理である。選択学習の結果は本書のモデルによる計算である。誤差を計算する回路は仮説である。

{\footnotesize
\setlength{\tabcolsep}{3pt}\renewcommand{\arraystretch}{1.15}
\begin{longtable}{>{\raggedright}p{0.5cm}>{\raggedright}p{4.0cm}>{\raggedright}p{5.6cm}>{\raggedright}p{2.3cm}>{\raggedright\arraybackslash}p{1.7cm}}
\toprule
No. & 主張 & 実験と測定内容 & 出典 & 状態 \\
\midrule\endfirsthead
\toprule
No. & 主張 & 実験と測定内容 & 出典 & 状態 \\
\midrule\endhead
1 & 人工ネットワークにおける形の誤差逆伝播は脳では見つかっていない & 直接の実験はない。これは不在についての主張である。レビューは、このアルゴリズムが何を要求するか（順方向をなぞる逆向きの結合、独立したフェーズ）と、どんな生物学的近似が提案されているかを検討している。 & \cite{Lillicrap2020, WhittingtonBogacz2019} & 仮説 \\
2 & 重みは放出部位の数、放出確率、1パケットへの応答に分解される：$w=N_s\,p\,A_1$~\eqref{eq:quantal} & 放出を低下させたカエルの神経筋シナプス。受け手の応答は、1パケットに対する自発的な微小応答の整数倍の段階で変動する。スパイクあたりのパケット数はポアソン則に従う。 & \cite{delCastillo1954} & 測定済み \\
3 & 受け手の受容体は同時検出器であり、カルシウムが重みの変化を起動する & マウスニューロンのパッチクランプ。静止電位ではMg$^{2+}$イオンがグルタミン酸で開いたチャネルを塞ぎ、脱分極で外れる。海馬スライス。受け手内のカルシウムキレート剤は長期増強を阻止し、細胞内で光によりカルシウムを放出すると、それだけで伝達が増強される。 & \cite{Nowak1984, Malenka1988calcium} & 測定済み \\
4 & 決定はどちらの側でも実行される：受け手で$A_1$が増え、送り手で$p$が変わる & 1個のスパインの上で光により放出したグルタミン酸は、スパインとその受容体を通る電流の速く持続的な増大を起こす。増強にはAMPA受容体の挿入が伴う。受け手の脱分極は送り手の放出を一時的に抑える。その効果は間隙を逆向きに進むエンドカンナビノイドが運ぶ（\nt{GABA}入力で示された）。短期的な枯渇は送り手の放出確率に依存する。 & \cite{Matsuzaki2004, Malinow2002, WilsonNicoll2001, Tsodyks1997} & 測定済み \\
5 & 送り手と受け手が同時に活動するとシナプスは強まる~\eqref{eq:hebb} & 麻酔下のウサギ。海馬の入力路への短いスパイク列の後、応答は$30$~分から$10$~時間増強される。海馬スライスでは、受け手のスパイクは同じ瞬間に活動しているシナプスだけを強める。 & \cite{Bliss1973LTP, Kelso1986hebb} & 測定済み \\
6 & 規則~\eqref{eq:oja}は入力の相関の主固有ベクトルを見つける（命題~\ref{prop:oja}） & 定理。証明は本章。ニューロンが入力の主成分を学んだという実験はない。シナプスで重みの増大を制限する機構は確立していない。 & \cite{Oja1982} & 理論 \\
7 & タイミングによるヘッブ則：窓は約$\pm20\ms$（図~\ref{fig:stdp}）。繰り返される系列から連鎖が育つ & 培養した海馬ニューロンのペア。送り手のスパイクの後$20\ms$以内の受け手のスパイクは持続的な増強を、前$20\ms$以内のものは抑圧を生み、どちらもNMDA受容体に依存する。図の比$A_-=0.6A_+$は説明用である。ネットワークでこうして連鎖が育つことは直接測定されていない。 & \cite{BiPoo1998} & 測定済み \\
8 & トレース~\eqref{eq:threefactor}：同時の活動の後$1$--$2\,\text{s}$に届いたドーパミンは結合を強め、それより遅いものは強めない（命題~\ref{prop:threefactor}） & 線条体スライスで、グルタミン酸入力とドーパミン入力を別々に光遺伝学的に刺激。ドーパミンがグルタミン酸入力の$0.3$--$2\,\text{s}$後に来たときだけスパインが大きくなる。窓はcAMPの速い上昇と分解で決まる。皮質では、ペアリングが増強と抑圧の別々のトレースを残し、モノアミン受容体がそれを秒程度の窓で長期的な変化に変える。 & \cite{Yagishita2014, He2015eligibility} & 測定済み \\
9 & ドーパミンなしでも秒単位の窓がある：第三の信号は受け手の強い興奮 & 生きたマウスのCA1野。受け手の1回のプラトー事象が、その前後数秒に来た入力を強め、1回の通過で場所野を作る。 & \cite{Bittner2017} & 測定済み \\
10 & 三要素則の平均ステップは報酬の勾配に沿う~\eqref{eq:perturbation}（命題~\ref{prop:gradient}） & ランダムな摂動による学習の勾配定理。本章では小さな雑音について導いた。 & \cite{Williams1992} & 理論 \\
11 & 雑音は探索である：ネットワークは自分のランダムなずれから学ぶ & 若いキンカチョウ。大脳基底核回路から出る核を不活性化すると、歌の変動性が急激かつ可逆的に消える。成鳥のキンカチョウ。平均より一方の側の高さで歌われた音節に妨害を与えると、鳥はすばやく高さを反対側にずらす。自分のばらつきから学ぶ。 & \cite{Olveczky2005, TumerBrainard2007} & 測定済み \\
12 & 二者択一は$\tau_e=1\,\text{s}$、$\delta=R-\bar R$で学習され、$\tau_e=50\ms$では学習されない。期待値を引かないと重みは際限なく増え、学習率が大きいとネットワークは悪い選択を固定しうる & \texttt{models/dofa\_learning.py}、$\eta=0.05$、$500$試行、シード$6$個。$\tau_e=1\,\text{s}$：$A$を選ぶ割合$0.65$--$0.79\to0.96$--$1.00$、重み$0.85$--$1.33$。$\tau_e=50\ms$：$0.38$--$0.55$、重みは変わらない。$\delta=R$：勝者の重み$860$--$1190$。$\delta=R$、$\eta=0.5$、シード$3$個：1個が$B$に固定（$0.04\to0$）。スクリプトは四つの場合をすべて出力する。再計算は本章と一致した。 & \texttt{models/\allowbreak dofa\_\allowbreak learning.py} & モデル \\
13 & 共通の信号一つによる学習時間は、雑音を出すニューロンの数に比例して伸びる（命題~\ref{prop:broadcastcost}） & 線形ネットワークの学習曲線。ニューロン摂動は正確な勾配より出力ニューロン数倍遅く、重み摂動はさらに入力数倍遅い。 & \cite{Werfel2005} & 理論 \\
14 & \nt{DOPA}の出力は主に行動選択の領域に向かう & ラットの黒質線条体\nt{DOPA}ニューロン個々の軸索の完全再構成。1細胞の枝は線条体体積の$0.45$--$5.7\,\%$（平均$2.7\,\%$）を覆い、線条体の外にはほとんど枝がない。 & \cite{Matsuda2009} & 測定済み \\
15 & 皮質は自分の入力を予測することを学び、誤差はその場で計算される & レビュー。感覚皮質には予期した入力と実際の入力の不一致に対する応答がある。これが皮質の一般的な学習則であることは証明されていない。 & \cite{Keller2018} & 仮説 \\
16 & ドーパミンは誤差~\eqref{eq:ctd}のようにふるまう：バースト、予告信号への移行、落ち込み（図~\ref{fig:ctdcases}） & サルの\nt{DOPA}ニューロン。予期しないジュースにバースト。学習後は予告信号にバーストが出て、約束のジュースには応答しない。ジュースが来ないと落ち込む。連続形~\eqref{eq:ctd}は理論である。 & \cite{Schultz1997, Doya2000} & 測定済み \\
17 & $\dd V/\dd t$を計算し期待を差し引く回路 & マウスの腹側被蓋野での記録と光遺伝学。\nt{DOPA}ニューロンは報酬から期待を差し引く。近くの\nt{GABA}ニューロンは報酬が期待されるときにそれを抑制し、その興奮や抑制は誤差を変える。微分がどう計算されるかは示されていない。 & \cite{Eshel2015arithmetic} & 仮説 \\
18 & \nt{DOPA}ニューロンはペースメーカーで、落ち込みはバーストより浅い & スライスで\nt{DOPA}ニューロンは、ゆっくりしたペースメーカー電位の上で自発的に非常に規則的に発火する。サルではレートは正の誤差を定量的に追うが、負の誤差は弱くしか伝えない。 & \cite{GraceOnn1989, Bayer2005} & 測定済み \\
19 & アクターとクリティック（図~\ref{fig:actorcritic}） & アルゴリズムは強化学習の理論から。ヒト（fMRI）では、腹側線条体は選択があってもなくてもクリティックのようにふるまい、背側線条体は行動を選ぶ必要があるときだけそうふるまう。 & \cite{SuttonBarto2018, ODoherty2004actor} & 仮説 \\
20 & 第二の受容体ファミリーをもつニューロンでは規則の$\delta$の符号が反転する & 線条体スライス。D1受容体をもつニューロンではペアリングがD1を要する増強を、D2をもつニューロンではD2を要する抑圧を生む。 & \cite{Shen2008} & 測定済み \\
\bottomrule
\end{longtable}}
